% =====================================================================
%  Chapter 1: Introduction and background        Page budget: 3 (pp. 1-3)
%  Booklet (ch. 3): a funnel - global trends -> application context ->
%           motivation (why existing approaches are not enough) ->
%           structure of the report.
%  Owner: Raneem (Claude Code) - only she edits this file
% =====================================================================
\chapter{מבוא ורקע}
\label{ch:intro}

דוח זה מציג תוכנית השקה בת \ltr{6} חודשים בישראל עבור \en{AURA}, מותג טינט
שפתיים בהשראת מזון שפותח בתרגיל \ltr{1}. קהל היעד הוא נשים בנות \ltr{18--28},
וערוצי ההשקה הם \en{TikTok} ו-\en{Instagram}, משפיעניות קטנות
(\en{Micro-Influencers}) ודוכני טעימות בקמפוסים. המבוא בנוי כמשפך: הוא פותח
בשלוש מגמות עולמיות, עובר להקשר הישראלי, מסביר מדוע דרכי ההשקה המקובלות אינן
מספיקות למותג חדש, ומסתיים במפת הדוח.

\section{מגמות עולמיות: ביוטי בהשראת מזון ומסחר חברתי}
\label{sec:intro-trends}

\textbf{ביוטי בהשראת מזון.} קטגוריית מוצרי השפתיים היא כיום מנוע צמיחה מרכזי
באיפור: בשוק הביוטי היוקרתי בארה"ב צמחה הקטגוריה ב-\ltr{19\%} בשנת \ltr{2024},
לעומת \ltr{5\%} באיפור כולו~\cite{circana2025beauty}, ומגמה זו נמשכה
ב-\ltr{2025}~\cite{glossy2025lips}. בתוך הקטגוריה מתחזק גל המוצרים בהשראת מזון
(\en{Food-Inspired Beauty}), שעבר משלב של גימיק לקטגוריה בפני
עצמה~\cite{waldeck2026foodbased}. דוגמה בולטת היא שפתון בהשראת חמוצים של
\en{e.l.f. Cosmetics}, שאזל תוך \ltr{8} דקות; לפי החברה, \ltr{77\%} מרוכשיו
היו לקוחות חדשים של המותג~\cite{waldeck2026foodbased}. גם חיפושי ניחוח ``גויאבה''
זינקו ב-\ltr{1,000\%} בשנה אחת, לפי נתוני \en{Spate}~\cite{jensen2025guava}.
אנשי התעשייה מסבירים את המשיכה בכך שמזון מוכר ונגיש לכל אחד: שם של פרי או
קינוח מעורר תחושה ואווירה שגוון לבדו אינו מעורר~\cite{jensen2025guava,waldeck2026foodbased}.

\textbf{מסחר חברתי ודור \en{Z}.} הרשתות החברתיות הפכו לנקודת הגילוי העיקרית של
מוצרי יופי בקרב דור \en{Z}. לפי סקר של סוכנות המשפיענים \en{Kyra},
\en{TikTok} מוביל כפלטפורמת גילוי (\ltr{39\%}), ואחריו \en{Instagram}
(\ltr{35\%}); נתח הרכישות דרך מסחר חברתי (\en{Social Commerce}) הוכפל
ל-\ltr{12\%}, ודור \en{Z} נוטה לנסות מותגים חדשים יותר מאשר לשמור נאמנות
למותג קיים~\cite{meintjes2025genz}. במקביל צומח שוק שיווק המשפיענים
(\en{Influencer Marketing}) מ-\ltr{\$21.1} מיליארד בשנת \ltr{2024}
ל-\ltr{\$24} מיליארד בשנת \ltr{2025}, ורוב היוצרים ב-\en{Instagram} (\ltr{76\%})
הם משפיעני ננו עם אלפי עוקבים בלבד~\cite{hypeauditor2025state}.

\section{ההקשר הישראלי}
\label{sec:intro-israel}

שוק מוצרי הקוסמטיקה בישראל מוערך ב-\ltr{\$408.53} מיליון בשנת \ltr{2026}, וצפוי
לצמוח עד \ltr{2031} בשיעור שנתי ממוצע (\en{CAGR}) של \ltr{4.60\%}. בין מנועי
הצמיחה נמנים מגזר היוקרה הנגישה (\en{Masstige}) והמסחר
המקוון~\cite{mordor2026israel}. התשתית הדיגיטלית לפנייה ישירה לקהל קיימת: בישראל
כ-\ltr{5.00} מיליון משתמשי \en{Instagram} וכ-\ltr{4.49} מיליון משתמשי
\en{TikTok} בגיל \ltr{18} ומעלה~\cite{datareportal2026israel}. קהל היעד מרוכז גם
פיזית: בשנת הלימודים תשפ"ז צפויים ללמוד כ-\ltr{350,900} סטודנטים ב-\ltr{57}
מוסדות~\cite{che2026students}, ונשים הן \ltr{60.3\%} מהלומדים לתואר
אקדמי~\cite{cbs2024students}. כפי שמפורט בפרק~\ref{ch:market}, מדובר בכ-\ltr{211}
אלף סטודנטיות --- קהל שאפשר להגיע אליו גם ברשת וגם בקמפוס.

עם זאת, השוק המקומי צפוף ומקוטב. מותגי הטינט הוויראליים מחו"ל נרכשים בישראל בעיקר דרך
יבוא מקביל, דיוטי פרי או משלוח מחו"ל, ומחירם לצרכנית מגיע ל-\ltr{99--175}
ש"ח~\cite{rhode2026terms,makeupstore2026rhode,caretobeauty2026benefit}; בקצה השני
עומדים מוצרי פארם זולים, סביב \ltr{15} ש"ח, ללא סיפור מותג~\cite{essence2026superpharm}.
בין שני הקצוות נותר פער מחיר, ש-\en{AURA} מכוונת אליו במחיר של \ltr{69} ש"ח
(פרק~\ref{ch:market}).

\section{מוטיבציה: מדוע הגישות הקיימות אינן מספיקות}
\label{sec:intro-motivation}

הדרכים המקובלות להשקת מותג קוסמטיקה אינן מתאימות למותג חדש וקטן בשוק הזה, משלוש
סיבות:

\begin{itemize}
  \item \textbf{מדף בפארם בלבד.} מוצר שעומד על מדף, לצד מותגים מוכרים, אינו
    בונה מודעות בעצמו. מותג שאיש לא שמע עליו זקוק לסיבה להיבחר.
  \item \textbf{פרסום ממומן או סלבריטאים.} תקציב השקה מוגבל אינו מאפשר קמפיין
    רחב, ודור \en{Z} מושפע מיוצרי תוכן לפחות כמו מחברים~\cite{meintjes2025genz}.
    משפיעניות קטנות זמינות בשפע: ברשת \en{Instagram}, \ltr{76\%} מכלל
    היוצרים הם משפיעני ננו~\cite{hypeauditor2025state}.
  \item \textbf{מכירה מקוונת בלבד.} מוצר שכל ערכו בגוון, במרקם ובניחוח קשה למכור
    מבלי לנסות אותו. בסקר של \ltr{20,000} צרכנים, \ltr{75\%} השיבו כי סביר
    יותר שירכשו ממותג שלא הכירו לאחר שקיבלו דוגמית, ורק \ltr{33\%} היו מגיעים
    לחנות כדי לנסותו~\cite{cosmeticsbusiness2021sampling}.
\end{itemize}

מכאן הרעיון שבבסיס התוכנית: לשלב תוכן ב-\en{TikTok} וב-\en{Instagram} עם
התנסות פיזית בדוכני טעימות בקמפוס, כך שכל ערוץ מכסה את החולשה של
האחר. פרק~\ref{ch:problem} מנסח זאת כבעיה וכיעדים מדידים.

\section{מבנה הדוח}
\label{sec:intro-structure}

הדוח בנוי כטיעון אחד רציף (איור~\ref{fig:roadmap}). פרק~\ref{ch:problem} מגדיר
את הבעיה ואת יעדי ההשקה; פרק~\ref{ch:literature} סוקר את הספרות על שיווק משפיענים,
מסחר חברתי ושיווק חושי, ומציג מאמר דגל; פרק~\ref{ch:market} אומד את היקף השוק
ומשווה מתחרים; פרק~\ref{ch:alternatives} משווה חלופות השקה ובוחר אחת מהן;
פרקים~\ref{ch:workplan} ו-\ref{ch:risks} מפרטים את תוכנית העבודה ואת ניהול
הסיכונים; ופרק~\ref{ch:summary} מסכם.

\begin{insightbox}[בפרויקט שלנו]
\textbf{המגמה:} ביוטי בהשראת מזון ומסחר חברתי. \textbf{ההקשר:} מותג חדש בשוק
קוסמטיקה ישראלי צפוף, עם כ-\ltr{211} אלף סטודנטיות פעילות ברשת.
\textbf{המוטיבציה:} מדף, פרסום רחב או חנות מקוונת לבדם אינם בונים מודעות ואינם
מאפשרים התנסות. \textbf{מטרת הדוח:} להראות שהשקה משולבת של רשתות חברתיות
ודוכני טעימות היא בת ביצוע בתוך \ltr{6} חודשים.
\end{insightbox}

\begin{figure}[htbp]
  \begin{diagram}
    % Report roadmap: chapters 1-8, read right to left.
% Used inside \begin{diagram} (LTR); Hebrew labels use \texthebrew.
% Colors follow the chapter shades of aura-report.cls.
\begin{tikzpicture}[
    step/.style={rectangle, rounded corners=5pt, draw=#1, line width=0.9pt,
      fill=#1!12, minimum width=1.55cm, minimum height=1.25cm,
      text width=1.5cm, align=center, font=\sffamily\footnotesize,
      text=AuraPlum},
    num/.style={circle, fill=#1, text=white, inner sep=1.4pt,
      font=\sffamily\scriptsize\bfseries},
    flow/.style={-{Stealth[length=5pt]}, AuraGold, line width=1pt}]
  \foreach \i/\c/\t in {
      1/AuraHoney/{מבוא\\ורקע},
      2/AuraPeach/{הבעיה\\והיעדים},
      3/AuraBerry/{ספרות\\ומאמר דגל},
      4/AuraRuby/{חקר\\שוק},
      5/AuraCocoa/{חלופות\\השקה},
      6/AuraHoney/{תוכנית\\עבודה},
      7/AuraPeach/{ניהול\\סיכונים},
      8/AuraBerry/{סיכום}} {
    \node[step=\c] (s\i) at ({-1.92*\i}, 0) {\texthebrew{\t}};
    \node[num=\c] at (s\i.north) {\i};
  }
  \foreach \i [evaluate=\i as \j using int(\i+1)] in {1,...,7}
    \draw[flow] (s\i.west) -- (s\j.east);
\end{tikzpicture}

  \end{diagram}
  \caption{מפת הדרכים של הדוח: כל פרק נשען על הפרק שלפניו}
  \label{fig:roadmap}
\end{figure}
